\documentclass[manuscript]{acmart}

\usepackage{array}
\usepackage{mathtools}
\usepackage{stmaryrd}
\usepackage{mathpartir}

\setcopyright{none}
\copyrightyear{}
\acmYear{}
\acmDOI{}
\acmConference{}{}{}
\acmISBN{}

\newcommand{\lang}{\textsc{SwiftyPOPL}}
\newcommand{\kw}[1]{\ensuremath{\mathsf{#1}}}
\newcommand{\Int}{\kw{Int}}
\newcommand{\Bool}{\kw{Bool}}
\newcommand{\String}{\kw{String}}
\newcommand{\Unit}{\kw{Unit}}
\newcommand{\Self}{\kw{Self}}
\newcommand{\Any}[1]{\kw{any}\ #1}
\newcommand{\ok}{\ \kw{ok}}
\newcommand{\dom}{\operatorname{dom}}
\newcommand{\fields}{\operatorname{fields}}
\newcommand{\methods}{\operatorname{methods}}
\newcommand{\reqs}{\operatorname{reqs}}
\newcommand{\params}{\operatorname{params}}
\newcommand{\ret}{\operatorname{ret}}
\newcommand{\labels}{\operatorname{labels}}
\newcommand{\subst}[3]{[#1/#2]#3}
\newcommand{\ftype}[3]{\forall #1.\ (#2) \to #3}
\newcommand{\emptyctx}{\cdot}
\newcommand{\has}{\vdash}
\newcommand{\modelsP}{\mathrel{\triangleright}}

\begin{document}

\title{A Surface Type System for \lang}

\author{Runbang Yan}
\email{rbyan25@stu.pku.edu.cn}
\affiliation{%
  \institution{School of Computer Science, Peking University}
  \city{Beijing}
  \country{China}
}

\renewcommand{\shortauthors}{Yan}

\begin{abstract}
This note gives a self-contained, proof-oriented type system for the surface language of \lang.
The system is intentionally syntax directed.
It separates nominal source concepts, such as structs, protocols, extensions, generic constraints, and \Self, from the lower-level representation used by a future elaboration pass.
The main meta-theoretic obligation of this surface system is elaboration preservation: every well-typed surface program elaborates to a well-typed intermediate representation, and that intermediate representation can then be lowered to a well-typed core term.
\end{abstract}

\maketitle

\section{Surface Syntax}

We write $S$ for struct names, $P,Q$ for protocol names, $X,Y$ for type variables, $x,f,m$ for value variables and function or method names, and $\ell$ for external argument labels.
The distinguished label $\epsilon$ denotes an unlabeled positional parameter.

\[
\begin{array}{rcll}
\tau & ::= & \Int \mid \Bool \mid \String \mid \Unit
  & \text{base types} \\
& \mid & S
  & \text{nominal struct type} \\
& \mid & X
  & \text{type variable} \\
& \mid & \Self
  & \text{contextual self type} \\
& \mid & \Any{P}
  & \text{existential protocol type} \\
& \mid & \ftype{\overline{X : \overline{P}}}{\overline{\ell : \tau}}{\tau}
  & \text{function scheme} \\[2mm]

e & ::= & n \mid b \mid s \mid x \mid \kw{self}
  & \text{literals and variables} \\
& \mid & \kw{func}\langle\overline{X : \overline{P}}\rangle
          (\overline{\ell\ x : \tau}) \to \tau\ \{ ss \}
  & \text{function value} \\
& \mid & e(\overline{a}) \mid e\langle\overline{\tau}\rangle
  & \text{call and type application} \\
& \mid & e.m \mid S(\overline{\ell : e})
  & \text{member access and construction} \\
& \mid & \kw{if}\ e\ \{ ss \}\ \kw{else}\ \{ ss \}
  & \text{conditional expression} \\
& \mid & e\ \kw{as}\ \Any{P}
  & \text{existential package} \\[2mm]

a & ::= & e \mid \ell : e
  & \text{positional or labeled argument} \\[2mm]

ss & ::= & \epsilon \mid stmt;\ ss
  & \text{statement sequence} \\
stmt & ::= & e \mid d
  & \text{expression or declaration} \\[2mm]

d & ::= & \kw{let}\ x : \tau = e
  & \text{value declaration} \\
& \mid & \kw{func}\ f\langle\overline{X : \overline{P}}\rangle
          (\overline{\ell\ x : \tau}) \to \tau\ \{ ss \}
  & \text{function definition} \\
& \mid & \kw{struct}\ S\ \{ \overline{\kw{let}\ x : \tau};\ \overline{md} \}
  & \text{struct definition} \\
& \mid & \kw{protocol}\ P\ \{ \overline{msig} \}
  & \text{protocol definition} \\
& \mid & \kw{extension}\ S : P\ \{ \overline{md} \}
  & \text{conformance definition} \\[2mm]

md & ::= & \kw{func}\ m\langle\overline{X : \overline{P}}\rangle
          (\overline{\ell\ x : \tau}) \to \tau\ \{ ss \}
  & \text{method definition} \\
msig & ::= & \kw{func}\ m\langle\overline{X : \overline{P}}\rangle
          (\overline{\ell : \tau}) \to \tau
  & \text{protocol method requirement}
\end{array}
\]

The implementation currently parses a subset of the grammar above.
The type system includes $\Any{P}$ because it is part of the intended surface language and is the source construct that elaborates to existential packages.

\section{Environments}

The surface type system uses four environments.

\[
\begin{array}{rcll}
\Psi & ::= &
  \left(
    \Psi_S,\Psi_P,\Psi_C,\Psi_M,\Psi_V
  \right)
  & \text{global environment} \\
\Psi_S(S) &=& \{\overline{x : \tau}\}
  & \text{struct fields} \\
\Psi_P(P) &=& \{\overline{m : \sigma}\}
  & \text{protocol requirements} \\
\Psi_C(S,P) &=& \{\overline{m : \sigma}\}
  & \text{extension methods implementing } S:P \\
\Psi_M(S) &=& \{\overline{m : \sigma}\}
  & \text{struct and extension methods of } S \\
\Psi_V(x) &=& \sigma
  & \text{top-level value or function type} \\[1mm]
\Omega & ::= & \emptyctx \mid \Omega, X : \overline{P}
  & \text{type-variable constraints} \\
\Gamma & ::= & \emptyctx \mid \Gamma, x : \tau
  & \text{local value environment} \\
\rho & ::= & \bot \mid S \mid P
  & \text{\Self resolver}
\end{array}
\]

A method scheme always contains an explicit first parameter for the receiver.
For a struct method on $S$, the first parameter has type $S$.
For a protocol requirement in $P$, the first parameter has type $\Self$ under self resolver $P$.
This convention matches the implementation's resolver pass and makes method calls reducible to ordinary function calls during elaboration.

\[
\sigma ::= \ftype{\overline{X : \overline{P}}}{\overline{\ell : \tau}}{\tau}
\]

We use the following auxiliary total functions in the rules.
They are finite-map lookups and list operations, so they can be defined by structural recursion.

\[
\begin{array}{rcl}
\params(\ftype{\overline{X : \overline{P}}}{\overline{\ell : \tau}}{\tau_r})
  &=& \overline{\ell : \tau} \\
\ret(\ftype{\overline{X : \overline{P}}}{\overline{\ell : \tau}}{\tau_r})
  &=& \tau_r \\
\labels(\overline{\ell : \tau}) &=& \overline{\ell}
\end{array}
\]

\section{Type Well-Formedness}

The judgment
\[
  \Psi;\Omega;\rho \has \tau \ok
\]
means that type $\tau$ is well formed under global environment $\Psi$, generic context $\Omega$, and current \Self resolver $\rho$.

\begin{mathpar}
\inferrule*[Right=\textsc{WF-Int}]
  { }
  { \Psi;\Omega;\rho \has \Int \ok }

\inferrule*[Right=\textsc{WF-Bool}]
  { }
  { \Psi;\Omega;\rho \has \Bool \ok }

\inferrule*[Right=\textsc{WF-String}]
  { }
  { \Psi;\Omega;\rho \has \String \ok }

\inferrule*[Right=\textsc{WF-Unit}]
  { }
  { \Psi;\Omega;\rho \has \Unit \ok }

\inferrule*[Right=\textsc{WF-Struct}]
  { S \in \dom(\Psi_S) }
  { \Psi;\Omega;\rho \has S \ok }

\inferrule*[Right=\textsc{WF-TVar}]
  { X \in \dom(\Omega) }
  { \Psi;\Omega;\rho \has X \ok }

\inferrule*[Right=\textsc{WF-Self}]
  { \rho \neq \bot }
  { \Psi;\Omega;\rho \has \Self \ok }

\inferrule*[Right=\textsc{WF-Any}]
  { P \in \dom(\Psi_P) }
  { \Psi;\Omega;\rho \has \Any{P} \ok }

\inferrule*[Right=\textsc{WF-Fun}]
  {
    \overline{P_i \in \dom(\Psi_P)} \\
    \Omega' = \Omega,\overline{X_i : \overline{P_i}} \\
    \overline{\Psi;\Omega';\rho \has \tau_j \ok} \\
    \Psi;\Omega';\rho \has \tau_r \ok
  }
  {
    \Psi;\Omega;\rho \has
      \ftype{\overline{X_i : \overline{P_i}}}
             {\overline{\ell_j : \tau_j}}
             {\tau_r}
    \ok
  }
\end{mathpar}

\Self is resolved only after well-formedness succeeds.
The partial function $\mathsf{resolve}_{\rho}$ is defined as follows.

\[
\begin{array}{rcl}
\mathsf{resolve}_{S}(\Self) &=& S \\
\mathsf{resolve}_{P}(\Self) &=& \Self \\
\mathsf{resolve}_{\rho}(\ftype{\overline{X : \overline{P}}}{\overline{\ell:\tau}}{\tau_r})
  &=&
  \ftype{\overline{X : \overline{P}}}
         {\overline{\ell:\mathsf{resolve}_{\rho}(\tau)}}
         {\mathsf{resolve}_{\rho}(\tau_r)} \\
\mathsf{resolve}_{\rho}(\tau) &=& \tau
  \quad \text{for every other type constructor.}
\end{array}
\]

Inside protocol declarations, $\Self$ remains abstract.
Inside struct and extension declarations, $\Self$ resolves to the concrete struct type.

\section{Protocol Satisfaction}

The judgment
\[
  \Psi;\Omega \has \tau \modelsP P
\]
means that type $\tau$ is known to satisfy protocol $P$.

\begin{mathpar}
\inferrule*[Right=\textsc{Sat-Struct}]
  { (S,P) \in \dom(\Psi_C) }
  { \Psi;\Omega \has S \modelsP P }

\inferrule*[Right=\textsc{Sat-TVar}]
  { X : \overline{P} \in \Omega \\ P \in \overline{P} }
  { \Psi;\Omega \has X \modelsP P }
\end{mathpar}

Only concrete structs and constrained type variables satisfy protocols.
In particular, $\Any{P}$ does not satisfy $P$ by default.
Values of type $\Any{P}$ must be explicitly opened by the elaborated core program before their packed operations can be used.
This restriction is the standard existential discipline and is what makes the later proof modular.

\section{Argument Matching}

The source language follows Swift-style external labels.
Rather than duplicate label-ordering details inside every typing rule, we use an auxiliary judgment:

\[
  \Psi;\Omega;\Gamma;\rho \has
  \overline{a} \Leftarrow \overline{\ell:\tau}
  \Rightarrow \overline{e:\tau}
\]

It means that actual arguments $\overline{a}$ match formal parameters $\overline{\ell:\tau}$, producing the ordered argument expressions $\overline{e}$.
The judgment is deterministic.
It enforces that positional arguments match $\epsilon$ parameters in order, labeled arguments match by label, labels are not duplicated, and no positional argument appears after a labeled argument.

\begin{mathpar}
\inferrule*[Right=\textsc{Args-Empty}]
  { }
  { \Psi;\Omega;\Gamma;\rho \has
    \epsilon \Leftarrow \epsilon \Rightarrow \epsilon }

\inferrule*[Right=\textsc{Args-Pos}]
  {
    \Psi;\Omega;\Gamma;\rho \has e : \tau \\
    \Psi;\Omega;\Gamma;\rho \has
      \overline{a} \Leftarrow \overline{\ell:\tau'}
      \Rightarrow \overline{e':\tau'}
  }
  {
    \Psi;\Omega;\Gamma;\rho \has
      e,\overline{a}
      \Leftarrow
      \epsilon:\tau,\overline{\ell:\tau'}
      \Rightarrow
      e:\tau,\overline{e':\tau'}
  }

\inferrule*[Right=\textsc{Args-Named}]
  {
    \ell_k = \ell \\
    \ell \notin \labels(\overline{a}) \\
    \Psi;\Omega;\Gamma;\rho \has e : \tau_k \\
    \Psi;\Omega;\Gamma;\rho \has
      \overline{a}
      \Leftarrow
      \overline{\ell_i:\tau_i}_{i \neq k}
      \Rightarrow
      \overline{e_i:\tau_i}_{i \neq k}
  }
  {
    \Psi;\Omega;\Gamma;\rho \has
      \ell:e,\overline{a}
      \Leftarrow
      \overline{\ell_i:\tau_i}_{i=1}^{n}
      \Rightarrow
      \overline{e_i:\tau_i}_{i=1}^{n}
  }
\end{mathpar}

The side condition in \textsc{Args-Named} ranges only over named arguments.
A complete mechanization would define two mutually recursive judgments, one for the positional prefix and one for the named suffix; the compact rule above is the declarative presentation of the same deterministic relation.

\section{Expression Typing}

The main expression typing judgment is
\[
  \Psi;\Omega;\Gamma;\rho \has e : \tau.
\]

\begin{mathpar}
\inferrule*[Right=\textsc{T-Int}]
  { }
  { \Psi;\Omega;\Gamma;\rho \has n : \Int }

\inferrule*[Right=\textsc{T-Bool}]
  { }
  { \Psi;\Omega;\Gamma;\rho \has b : \Bool }

\inferrule*[Right=\textsc{T-String}]
  { }
  { \Psi;\Omega;\Gamma;\rho \has s : \String }

\inferrule*[Right=\textsc{T-Var-Local}]
  { x:\tau \in \Gamma }
  { \Psi;\Omega;\Gamma;\rho \has x : \tau }

\inferrule*[Right=\textsc{T-Var-Global}]
  { x:\tau \in \Psi_V }
  { \Psi;\Omega;\Gamma;\rho \has x : \tau }

\inferrule*[Right=\textsc{T-Self-Struct}]
  { \rho = S }
  { \Psi;\Omega;\Gamma;\rho \has \kw{self} : S }

\inferrule*[Right=\textsc{T-Self-Protocol}]
  { \rho = P }
  { \Psi;\Omega;\Gamma;\rho \has \kw{self} : \Self }

\inferrule*[Right=\textsc{T-Fun}]
  {
    \overline{P_i \in \dom(\Psi_P)} \\
    \Omega' = \Omega,\overline{X_i : \overline{P_i}} \\
    \overline{\Psi;\Omega';\rho \has \tau_j \ok} \\
    \Psi;\Omega';\rho \has \tau_r \ok \\
    \Psi;\Omega';\Gamma,\overline{x_j:\tau_j};\rho \has ss : \tau_r
  }
  {
    \Psi;\Omega;\Gamma;\rho \has
    \kw{func}\langle\overline{X_i : \overline{P_i}}\rangle
      (\overline{\ell_j\ x_j : \tau_j}) \to \tau_r\ \{ ss \}
    :
    \ftype{\overline{X_i : \overline{P_i}}}
           {\overline{\ell_j:\tau_j}}
           {\tau_r}
  }

\inferrule*[Right=\textsc{T-TypeApp}]
  {
    \Psi;\Omega;\Gamma;\rho \has e :
      \ftype{\overline{X_i : \overline{P_i}}}
             {\overline{\ell_j:\tau_j}}
             {\tau_r} \\
    |\overline{X_i}| = |\overline{\tau_i'}| \\
    \overline{\Psi;\Omega;\rho \has \tau_i' \ok} \\
    \overline{\Psi;\Omega \has \tau_i' \modelsP P \quad
      \text{for every } P \in \overline{P_i}}
  }
  {
    \Psi;\Omega;\Gamma;\rho \has
      e\langle\overline{\tau_i'}\rangle
    :
      \subst{\overline{\tau_i'}}{\overline{X_i}}
      \left(
        \ftype{\emptyctx}{\overline{\ell_j:\tau_j}}{\tau_r}
      \right)
  }

\inferrule*[Right=\textsc{T-Call}]
  {
    \Psi;\Omega;\Gamma;\rho \has e_f :
      \ftype{\emptyctx}{\overline{\ell_i:\tau_i}}{\tau_r} \\
    \Psi;\Omega;\Gamma;\rho \has
      \overline{a} \Leftarrow \overline{\ell_i:\tau_i}
      \Rightarrow \overline{e_i:\tau_i}
  }
  {
    \Psi;\Omega;\Gamma;\rho \has
      e_f(\overline{a}) : \tau_r
  }

\inferrule*[Right=\textsc{T-Construct}]
  {
    \Psi_S(S) = \{\overline{x_i:\tau_i}\} \\
    \Psi;\Omega;\Gamma;\rho \has
      \overline{\ell_i:e_i}
      \Leftarrow
      \overline{x_i:\tau_i}
      \Rightarrow
      \overline{e_i:\tau_i}
  }
  {
    \Psi;\Omega;\Gamma;\rho \has
      S(\overline{\ell_i:e_i}) : S
  }

\inferrule*[Right=\textsc{T-Field}]
  {
    \Psi;\Omega;\Gamma;\rho \has e : S \\
    \Psi_S(S)(x) = \tau
  }
  {
    \Psi;\Omega;\Gamma;\rho \has e.x : \tau
  }

\inferrule*[Right=\textsc{T-StructMethod}]
  {
    \Psi;\Omega;\Gamma;\rho \has e : S \\
    \Psi_M(S)(m) =
      \ftype{\overline{X_i:\overline{P_i}}}
             {\epsilon:S,\overline{\ell_j:\tau_j}}
             {\tau_r}
  }
  {
    \Psi;\Omega;\Gamma;\rho \has e.m :
      \ftype{\overline{X_i:\overline{P_i}}}
             {\overline{\ell_j:\tau_j}}
             {\tau_r}
  }

\inferrule*[Right=\textsc{T-ConstrainedMethod}]
  {
    \Psi;\Omega;\Gamma;\rho \has e : X \\
    X : \overline{P} \in \Omega \\
    P_k \in \overline{P} \\
    \Psi_P(P_k)(m) =
      \ftype{\overline{Y_i:\overline{Q_i}}}
             {\epsilon:\Self,\overline{\ell_j:\tau_j}}
             {\tau_r} \\
    m \text{ is unambiguous among } \overline{P}
  }
  {
    \Psi;\Omega;\Gamma;\rho \has e.m :
      \subst{X}{\Self}
      \left(
        \ftype{\overline{Y_i:\overline{Q_i}}}
               {\overline{\ell_j:\tau_j}}
               {\tau_r}
      \right)
  }

\inferrule*[Right=\textsc{T-IfElse}]
  {
    \Psi;\Omega;\Gamma;\rho \has e_c : \Bool \\
    \Psi;\Omega;\Gamma;\rho \has ss_1 : \tau \\
    \Psi;\Omega;\Gamma;\rho \has ss_2 : \tau
  }
  {
    \Psi;\Omega;\Gamma;\rho \has
      \kw{if}\ e_c\ \{ ss_1 \}\ \kw{else}\ \{ ss_2 \}
      : \tau
  }

\inferrule*[Right=\textsc{T-AsAny}]
  {
    \Psi;\Omega;\Gamma;\rho \has e : \tau \\
    \Psi;\Omega \has \tau \modelsP P
  }
  {
    \Psi;\Omega;\Gamma;\rho \has
      e\ \kw{as}\ \Any{P} : \Any{P}
  }
\end{mathpar}

The source language intentionally has no elimination form for $\Any{P}$.
Opening an existential is inserted by elaboration at controlled sites, or added as an explicit future source construct.
This keeps the surface type system small and prevents accidental treatment of an existential value as if it were a concrete conforming type.

\section{Statement Sequences}

The judgment
\[
  \Psi;\Omega;\Gamma;\rho \has ss : \tau
\]
assigns a type to a statement sequence.
If the sequence contains no expression statement, its type is $\Unit$.
If it contains expression statements, its type is the type of the last expression statement.
Declarations extend the local environment only for the following statements in the same sequence.

\begin{mathpar}
\inferrule*[Right=\textsc{S-Empty}]
  { }
  { \Psi;\Omega;\Gamma;\rho \has \epsilon : \Unit }

\inferrule*[Right=\textsc{S-Expr-Last}]
  {
    \Psi;\Omega;\Gamma;\rho \has e : \tau
  }
  {
    \Psi;\Omega;\Gamma;\rho \has e;\epsilon : \tau
  }

\inferrule*[Right=\textsc{S-Expr-Cons}]
  {
    \Psi;\Omega;\Gamma;\rho \has e : \tau_1 \\
    \Psi;\Omega;\Gamma;\rho \has ss : \tau_2
  }
  {
    \Psi;\Omega;\Gamma;\rho \has e;ss : \tau_2
  }

\inferrule*[Right=\textsc{S-Let}]
  {
    \Psi;\Omega;\Gamma;\rho \has e : \tau \\
    \Psi;\Omega;\rho \has \tau \ok \\
    \Psi;\Omega;\Gamma,x:\tau;\rho \has ss : \tau_r
  }
  {
    \Psi;\Omega;\Gamma;\rho \has
      \kw{let}\ x:\tau = e;\ ss : \tau_r
  }

\inferrule*[Right=\textsc{S-Fun}]
  {
    \sigma =
      \ftype{\overline{X_i:\overline{P_i}}}
             {\overline{\ell_j:\tau_j}}
             {\tau_r} \\
    \Omega' = \Omega,\overline{X_i:\overline{P_i}} \\
    \overline{P_i \subseteq \dom(\Psi_P)} \\
    \overline{\Psi;\Omega';\rho \has \tau_j \ok} \\
    \Psi;\Omega';\rho \has \tau_r \ok \\
    \Psi;\Omega';\Gamma,f:\sigma,\overline{x_j:\tau_j};\rho
      \has ss_f : \tau_r \\
    \Psi;\Omega;\Gamma,f:\sigma;\rho \has ss : \tau
  }
  {
    \Psi;\Omega;\Gamma;\rho \has
      \kw{func}\ f\langle\overline{X_i:\overline{P_i}}\rangle
      (\overline{\ell_j\ x_j:\tau_j}) \to \tau_r\ \{ ss_f \};\ ss
      : \tau
  }
\end{mathpar}

The \textsc{S-Fun} rule is recursive: the function name is in scope in its own body.
If recursive functions are not desired, remove $f:\sigma$ from the premise that checks the function body.
The current implementation does not provide general recursion at the surface level; the rule above is the conservative specification to choose if recursive top-level functions are admitted later.

\section{Method and Declaration Well-Formedness}

Declarations are checked against a global environment that has already been collected from the program's top-level names.
This mirrors the two-pass implementation:
first collect names and declared signatures, then check bodies.

\subsection{Method Bodies}

The method typing judgment is
\[
  \Psi;\Omega;S \has md : \sigma
\]
for struct or extension methods, and
\[
  \Psi;\Omega;P \has msig : \sigma
\]
for protocol requirements.

\begin{mathpar}
\inferrule*[Right=\textsc{M-StructDef}]
  {
    \Omega' = \Omega,\overline{X_i:\overline{P_i}} \\
    \overline{\Psi;\Omega';S \has \tau_j \ok} \\
    \Psi;\Omega';S \has \tau_r \ok \\
    \Psi;\Omega';
      \kw{self}:S,\overline{x_j:\tau_j};
      S
      \has ss : \tau_r
  }
  {
    \Psi;\Omega;S \has
    \kw{func}\ m\langle\overline{X_i:\overline{P_i}}\rangle
      (\overline{\ell_j\ x_j:\tau_j}) \to \tau_r\ \{ ss \}
    :
    \ftype{\overline{X_i:\overline{P_i}}}
           {\epsilon:S,\overline{\ell_j:\tau_j}}
           {\tau_r}
  }

\inferrule*[Right=\textsc{M-ProtocolSig}]
  {
    \Omega' = \Omega,\overline{X_i:\overline{P_i}} \\
    \overline{\Psi;\Omega';P \has \tau_j \ok} \\
    \Psi;\Omega';P \has \tau_r \ok
  }
  {
    \Psi;\Omega;P \has
    \kw{func}\ m\langle\overline{X_i:\overline{P_i}}\rangle
      (\overline{\ell_j:\tau_j}) \to \tau_r
    :
    \ftype{\overline{X_i:\overline{P_i}}}
           {\epsilon:\Self,\overline{\ell_j:\tau_j}}
           {\tau_r}
  }
\end{mathpar}

The expression $\Psi;\Omega;S \has \tau \ok$ abbreviates the well-formedness judgment $\Psi;\Omega;\rho \has \tau \ok$ with resolver $\rho=S$.
Likewise, $\Psi;\Omega;P \has \tau \ok$ uses resolver $\rho=P$.

\subsection{Conformance}

A concrete method implementation matches a protocol requirement after substituting the receiver struct for the protocol's abstract \Self.

\[
\mathsf{match}(S,P,\sigma_{\mathit{impl}},\sigma_{\mathit{req}})
\quad \text{iff} \quad
\sigma_{\mathit{impl}} =
\subst{S}{\Self}\sigma_{\mathit{req}}.
\]

\begin{mathpar}
\inferrule*[Right=\textsc{D-Struct}]
  {
    \Psi_S(S)=\{\overline{x_i:\tau_i}\} \\
    \overline{\Psi;\emptyctx;S \has \tau_i \ok} \\
    \overline{\Psi;\emptyctx;S \has md_j : \sigma_j} \\
    \Psi_M(S) \supseteq \{\overline{m_j:\sigma_j}\}
  }
  {
    \Psi \has
      \kw{struct}\ S\ \{
        \overline{\kw{let}\ x_i:\tau_i};\
        \overline{md_j}
      \}
      \ok
  }

\inferrule*[Right=\textsc{D-Protocol}]
  {
    \Psi_P(P)=\{\overline{m_i:\sigma_i}\} \\
    \overline{\Psi;\emptyctx;P \has msig_i : \sigma_i}
  }
  {
    \Psi \has
      \kw{protocol}\ P\ \{\overline{msig_i}\}
      \ok
  }

\inferrule*[Right=\textsc{D-Extension}]
  {
    S \in \dom(\Psi_S) \\
    P \in \dom(\Psi_P) \\
    \Psi_P(P)=\{\overline{m_i:\sigma_i}\} \\
    \overline{\Psi;\emptyctx;S \has md_i : \sigma_i'} \\
    \overline{\mathsf{match}(S,P,\sigma_i',\sigma_i)}
  }
  {
    \Psi \has
      \kw{extension}\ S:P\ \{\overline{md_i}\}
      \ok
  }

\inferrule*[Right=\textsc{D-Let}]
  {
    \Psi;\emptyctx;\emptyctx;\bot \has e : \tau \\
    \Psi;\emptyctx;\bot \has \tau \ok
  }
  {
    \Psi \has \kw{let}\ x:\tau = e \ok
  }

\inferrule*[Right=\textsc{D-Fun}]
  {
    \sigma =
      \ftype{\overline{X_i:\overline{P_i}}}
             {\overline{\ell_j:\tau_j}}
             {\tau_r} \\
    f:\sigma \in \Psi_V \\
    \Omega' = \overline{X_i:\overline{P_i}} \\
    \overline{P_i \subseteq \dom(\Psi_P)} \\
    \overline{\Psi;\Omega';\bot \has \tau_j \ok} \\
    \Psi;\Omega';\bot \has \tau_r \ok \\
    \Psi;\Omega';f:\sigma,\overline{x_j:\tau_j};\bot
      \has ss : \tau_r
  }
  {
    \Psi \has
      \kw{func}\ f\langle\overline{X_i:\overline{P_i}}\rangle
      (\overline{\ell_j\ x_j:\tau_j}) \to \tau_r\ \{ ss \}
      \ok
  }
\end{mathpar}

\section{Program Well-Formedness}

A program is well typed if its top-level declaration collection produces a well-formed global environment and every top-level declaration checks against that environment.

\[
\begin{array}{rcl}
\mathsf{collect}(\overline{d}) &=& \Psi
\end{array}
\]

The collection function is defined only when all top-level names are distinct, all referenced struct and protocol names are declared, and all method names inside a single struct, protocol, or extension are distinct.
It records:
\[
\begin{array}{rcl}
\Psi_S(S) &=& \text{declared fields of } S, \\
\Psi_P(P) &=& \text{declared method requirements of } P, \\
\Psi_C(S,P) &=& \text{method implementations in } \kw{extension}\ S:P, \\
\Psi_M(S) &=& \text{methods declared inside } S \text{ plus extension methods}, \\
\Psi_V(x) &=& \text{declared type of top-level value or function } x.
\end{array}
\]

\begin{mathpar}
\inferrule*[Right=\textsc{P-Program}]
  {
    \mathsf{collect}(\overline{d}) = \Psi \\
    \overline{\Psi \has d_i \ok}
  }
  {
    \has \overline{d_i} \ok
  }
\end{mathpar}

\section{Meta-Theoretic Shape}

This surface system is designed to support a standard proof pipeline.
The surface language itself need not have an operational semantics.
Instead, it elaborates to a typed intermediate language.

\paragraph{Weakening.}
If $\Psi;\Omega;\Gamma;\rho \has e:\tau$, then adding fresh bindings to $\Gamma$ or adding unused type variables to $\Omega$ preserves the judgment.

\paragraph{Type substitution.}
If $\Psi;\Omega,X:\overline{P},\Omega';\Gamma;\rho \has e:\tau$ and $\Psi;\Omega \has \tau' \modelsP P$ for every $P \in \overline{P}$, then
\[
  \Psi;\Omega,\subst{\tau'}{X}\Omega';
    \subst{\tau'}{X}\Gamma;\rho
  \has
  \subst{\tau'}{X}e :
  \subst{\tau'}{X}\tau.
\]

\paragraph{Value substitution.}
If $\Psi;\Omega;\Gamma,x:\tau_x;\rho \has e:\tau$ and
$\Psi;\Omega;\Gamma;\rho \has v:\tau_x$, then
\[
  \Psi;\Omega;\Gamma;\rho \has \subst{v}{x}e:\tau.
\]

\paragraph{Elaboration preservation.}
Let $\llbracket \cdot \rrbracket$ translate surface types and environments to the chosen intermediate representation.
If the elaboration judgment is written
\[
  \Psi;\Omega;\Gamma;\rho \has e:\tau \leadsto E,
\]
then the preservation theorem should be:
\[
  \text{if } \Psi;\Omega;\Gamma;\rho \has e:\tau \leadsto E,
  \text{ then }
  \llbracket\Psi;\Omega;\Gamma;\rho\rrbracket
  \has_{\mathit{IR}}
  E : \llbracket\tau\rrbracket.
\]

\paragraph{End-to-end type safety.}
If $\has p \ok$, $p \leadsto I$, and $I \Downarrow t$, then $t$ is well typed in the core language.
Core progress and preservation then imply that the elaborated program cannot get stuck.

\end{document}
